\section{Evaluation of the degree-minus-four image}
\label{sch:sec:evaluation}

The first higher scalar map has rank one. Its kernel has a geometric
generator obtained by decorating the internal complete flag. The
calculation begins with the rank-two coefficient carried by that decoration.

Let $\nu_2$ be the complete-flag coefficient map. On its scalar
fibre $\mathbb P^1$, let $h=c_1(\mathcal O_{\mathbb P^1}(1))$,
with $\int_{\mathbb P^1}h=1$. Define
\[
 w_2=H^{-2}(\nu_2)(h)\in H^{-2}(P_{2,0}).
\]
The input shift of $\nu_2$ is $[4]$, so this notation uses degree
two on the flag fibre.

\begin{lemma}[The decorated rank-two flag]
\label{sch:lem:rank-two-decoration}
The class $w_2$ is nonzero. Thus it generates
$H^{-2}(P_{2,0})=\mathbb Q$.
\end{lemma}

\begin{proof}
In traceless matrix coordinates the potential is $2\det T$ on
$W=\operatorname{Mat}_3(\mathbb C)$. Scalar matrix coordinates
are free. Write $D\subset W$ for the determinant hypersurface.
Let $G_0=\operatorname{Gr}(2,3)=\mathbb P^2$, with tautological
plane bundle $U$, and let
\[
 \widetilde D=\operatorname{Hom}(\mathbb C^3,U)
      \longrightarrow D
\]
be the proper incidence remembering a plane containing the columns
of $T$. Its total space is smooth of complex dimension eight.
Write $z=c_1(\mathcal O_{G_0}(1))$. The tautological exact sequence
gives $c_1(U)=-z$ and $c_2(U)=z^2$.

Let $D'\subset D$ be the matrices of rank at most one and
$E'\subset\widetilde D$ its inverse image. Their complex dimensions
are five and six. Over $D\setminus D'$, the incidence is an
isomorphism. The proper resolution of $E'$ that remembers the
image line has total space
\[
 Y=\operatorname{Hom}(\mathbb C^3,L)
       \longrightarrow\operatorname{Fl}(1,2;3),
\]
where $L\subset U$ is the tautological line.
It is a rank-three vector bundle over the complete flag variety.
Its map to $E'$ is an isomorphism away from zero matrices.
The exceptional locus has complex dimension three, and its image
has complex dimension two. Closed-open exact sequences therefore
identify $H^{\mathrm{BM}}_{11}(E')$ with
$H^{\mathrm{BM}}_{11}(Y)=0$.
The last equality follows from the even cohomology of the complete
flag variety and the vector-bundle Thom isomorphism.

The top class of $E'$ maps to zero in
$H^{\mathrm{BM}}_{12}(\widetilde D)$.
To compute it, let
$p:\operatorname{Fl}(1,2;3)\to G_0$ forget the line and put
$q=c_1(U/L)$. The inclusion of $Y$ into the pullback of
$\widetilde D$ has quotient normal bundle
$\operatorname{Hom}(\mathbb C^3,U/L)$.
Its Euler class is $q^3$. Hence the class of $E'$ in the smooth
space $\widetilde D$ is Poincar\'e dual to
\[
 p_*(q^3)=c_1(U)^2-c_2(U)=z^2-z^2=0.
\]
Indeed, $q^2-c_1(U)q+c_2(U)=0$ by the tautological exact sequence,
while integration on the projective-line fibre gives $p_*q=1$
and $p_*1=0$. The resolution $Y\to E'$ is generically one-to-one,
so it pushes its fundamental class to that of $E'$.

The closed-open sequence for $E'\subset\widetilde D$ now shows
that
$H^{\mathrm{BM}}_{12}(\widetilde D)\to
 H^{\mathrm{BM}}_{12}(\widetilde D\setminus E')$
is an isomorphism. Its kernel is the zero exceptional class and
its cokernel maps to $H^{\mathrm{BM}}_{11}(E')=0$.
Since $\dim_{\mathbb C}D'=5$, the corresponding map from
$H^{\mathrm{BM}}_{12}(D)$ to the open stratum is also an isomorphism.
Consequently the proper pushforward is an isomorphism
\[
 H^{\mathrm{BM}}_{12}(\widetilde D)=\mathbb Q
       \xrightarrow{\sim}H^{\mathrm{BM}}_{12}(D)=\mathbb Q.
\]

A line in $\mathbb C^2$ determines its traceless stabilizer plane
in $\mathfrak{sl}_2$. Under the trace pairing, the perpendicular
line consists of nilpotent matrices with that image and kernel.
Thus the map from the original flag $\mathbb P^1$ to $G_0$
is the smooth conic $j:\mathbb P^1\hookrightarrow\mathbb P^2$.
Complex-oriented integration gives
\[
 j_*1=2z,\qquad j_*h=z^2.
\]
The complete-flag incidence is the restriction of $\widetilde D$
to this conic. Transitivity of complex Thom maps and proper
integration identifies its Gysin map with $j_*$ followed by the
universal determinant incidence map.

The support form of this universal map has the following description.
Poincar\'e duality on $\widetilde D$ takes $z^2$ to its nonzero
degree-twelve Borel--Moore class. Proper pushforward takes this to
the nonzero degree-twelve class of $D$.
The closed-open sequence in $W$, of real dimension eighteen,
takes it to the nonzero class in $H^5(\operatorname{GL}_3(\mathbb C))$.
The support inclusion
$D\subset\{\operatorname{Re}(2\det T)\leq0\}$ restricts this
class to the determinant Milnor fibre.
These are the maps of the same relative-cochain sequences defining
the normalized coefficient map, so the description also fixes its sign.

The determinant fibration has a product trivialization
$\operatorname{GL}_3(\mathbb C)\simeq
 \operatorname{SL}_3(\mathbb C)\times\mathbb C^*$.
Since $H^4(\operatorname{SL}_3)=0$, restriction induces an isomorphism
on $H^5$. Multiplication by the positive constant two in the
potential preserves the chosen support half-plane.
Therefore the image of $j_*h=z^2$ is nonzero in reduced Milnor
degree five, hence in normalized scalar degree $-2$.
This image is exactly $w_2$.
\end{proof}

The rank-three target in the required degree also has dimension one.
The two possible eigenvalue collisions determine this dimension.

\begin{lemma}[The next Borel--Moore group of the commuting-pair variety]
\label{sch:lem:rank-three-bm}
For the variety $\mathcal C_3$ of commuting \emph{pairs} of
three-by-three matrices,
\[
 H^{\mathrm{BM}}_{22}(\mathcal C_3;\mathbb Q)=\mathbb Q,
 \qquad H^{-4}(P_{3,0})=\mathbb Q.
\]
\end{lemma}

\begin{proof}
Remove the locus where $A$ has at most one eigenvalue.
The Jordan dimension calculation gives complex dimension at most
ten for this locus. Its removal does not change Borel--Moore homology
in degrees twenty-two and twenty-three. Denote the remaining space
by $X$, and let $O\subset X$ be the open set where $A$ has three
distinct eigenvalues.

The smooth variety $O$ has complex dimension twelve. Up to an
equivariantly contractible affine factor, its ordered eigenvalue cover is
\[
 \operatorname{GL}_3(\mathbb C)/T\times
                  \operatorname{Conf}_3(\mathbb C),
\]
where $T$ is the diagonal torus and $\operatorname{Conf}_3$ is
the space of ordered distinct triples. The symmetric group permutes
both factors, freely on the second. Rational cohomology of $O$
is the invariant part of this cover.
The first factor has the cohomology of the complete flag variety.
Its $H^1$ is zero, and its $H^2$ is generated by three Chern roots
with their sum zero. The permutation action has no invariant in
this degree.

For the second factor, write
\[
 (z_1,z_2,z_3)=(a,a+b,a+bc),\qquad
 a\in\mathbb C,\ b\in\mathbb C^*,\
 c\in\mathbb C\setminus\{0,1\}.
\]
These are global coordinates. Its degree-two cohomology has basis
\[
 \alpha=d\log b\wedge d\log c,\qquad
 \beta=d\log b\wedge d\log(c-1),
\]
where each logarithmic degree-one class includes the factor $(2\pi i)^{-1}$.
The transposition $(12)$ exchanges $\alpha$ and $\beta$.
The transposition $(23)$ sends
$\alpha\mapsto-\alpha$ and $\beta\mapsto\beta-\alpha$.
The term $d\log c\wedge d\log(c-1)$ has zero cohomology class
because it comes from the punctured $c$-line.
These two transformations have no common fixed vector.
In degree one, the three classes $d\log(z_i-z_j)$ are permuted,
so their invariant subspace has dimension one.
It follows that
\[
 H^1(O;\mathbb Q)=\mathbb Q,\qquad H^2(O;\mathbb Q)=0.
\]
Poincar\'e duality gives
$H^{\mathrm{BM}}_{23}(O)=\mathbb Q$ and
$H^{\mathrm{BM}}_{22}(O)=0$.

The closed complement $D_0=X\setminus O$ has two irreducible
components of complex dimension eleven. On the first, the repeated
eigenvalue of $A$ has a Jordan block of size two. On the second,
its two-dimensional eigenspace is semisimple.
Both strata are irreducible: the centralizer vector bundle over
each Jordan stratum gives their parameterization.
Their closures are distinct. The generic commuting matrix in the
semisimple stratum has two different eigenvalues on that eigenspace.
On the Jordan stratum the corresponding restriction is a polynomial
in the Jordan block, so its two eigenvalues coincide even in a limit.
The opposite containment is excluded by the smaller minimal polynomial
on the semisimple stratum. Thus
$H^{\mathrm{BM}}_{22}(D_0)=\mathbb Q^2$.

The boundary map from $H^{\mathrm{BM}}_{23}(O)$ to these two
fundamental classes is nonzero. More precisely, the discriminant
of the characteristic polynomial has multiplicities one and two
on the two components. At a generic Jordan point, $A$ is cyclic,
so the commuting pair variety is locally a centralizer vector bundle
over matrix space. The family
\[
 A_t=\begin{pmatrix}0&1&0\\t&0&0\\0&0&2\end{pmatrix}
\]
has characteristic discriminant $4t(4-t)^2$, of order one.
At a generic semisimple point, choose the commuting matrix $C$
with three distinct eigenvalues. The local commuting-pair chart then
diagonalizes $A$ in its eigenbasis. The family
\[
 A_t=\operatorname{diag}(t,-t,2)
\]
has discriminant $4t^2(4-t^2)^2$, of order two.
Both are transverse local coordinates in the specified smooth charts.
The linking class $d\log\operatorname{disc}(A)$ therefore has
boundary a nonzero common multiple of $(1,2)$, with complex orientations.
The closed-open sequence gives
\[
 H^{\mathrm{BM}}_{22}(X)
       =\mathbb Q^2/\mathbb Q(1,2)=\mathbb Q.
\]
Restoring the removed locus preserves this group.
Proposition~\ref{sch:prop:scalar} identifies it with
$H^{-4}(P_{3,0})$.
\end{proof}

The complete flag in dimension three has successive line bundles
$L_1,L_2,L_3$. Let $\ell_i=c_1(L_i)$ and define the three
degree-zero morphisms
\[
 \kappa_i:\mathbb Q_{\mathfrak M_3}[4]
   \longrightarrow Rb_{3*}\mathbb Q_{\mathcal T_3}[4]
   \xrightarrow{\ \ell_i\ }Rb_{3*}\mathbb Q_{\mathcal T_3}[6]
   \xrightarrow{\nu_3}\mathcal P_3.
\]
The first map is the adjunction unit. Cup product by a first Chern
class has cohomological degree two.

\begin{lemma}[Three line decorations]
\label{sch:lem:three-decorations}
At the scalar atlas point,
\[
 \kappa_3(1)\ne0,\qquad \kappa_2(1)=0
                    \quad\text{in }H^{-4}(P_{3,0}).
\]
At the critical point
$x=(\operatorname{diag}(0,0,1),0,0)$, identify the degree-$-4$
coefficient with the rank-two class $w_2$ tensored with the
normalized line class. Then
\[
 \kappa_1(1)=-2w_2,\qquad \kappa_2(1)=0,
 \qquad \kappa_3(1)=2w_2.
\]
\end{lemma}

\begin{proof}
The spectral projector for the two separated eigenvalue clusters
gives a local analytic decomposition into a plane and a line.
In this chart, $A$ is block diagonal. Its cross-block commutator
operator is invertible. The expression $\operatorname{Tr}(B[C,A])$
therefore splits exactly into the rank-two potential and a
nondegenerate bilinear form in the cross-block $B,C$ coordinates.
Rescaling the cross-block $C$ coordinates gives two eigenline-pair
quadratics. The rescaling is complex linear, so it preserves orientation.
Their normalized coefficient, including the free line variables,
identifies the target coefficient with $\mathcal P_2\boxtimes\mathcal P_1$.

An invariant complete flag separates according to the position of
the isolated eigenline. Its fibre at $x$ has three components,
each a copy of $\mathbb P^1$. On each component, the remaining
flag is the complete flag in the double eigenspace.
The two cross eigenline pairs contribute $(-1)^2=1$ in the
$B$ polarization. The conversion to the specified generators also
has ratio $(-1)^{p_2+p_1-p_3}=1$.
The flag tangents and lower $A$ normals cancel complex linearly.
Thus each component contributes the actual map $\nu_2$, with
coefficient $+1$, to the target identification just given.

Let $h$ be the positive hyperplane class on that component.
The tautological subline in the double eigenspace has Chern class
$-h$, and its quotient line has Chern class $h$.
The isolated line is constant. Accordingly the three restrictions are
\[
\begin{array}{c|ccc}
 \text{position of isolated line}&\ell_1&\ell_2&\ell_3\\\hline
 1&0&-h&h\\
 2&-h&0&h\\
 3&-h&h&0.
\end{array}
\]
Applying $\nu_2(h)=w_2$ and adding the three components proves
the values at $x$.

Each $\kappa_i$ induces a section of
$\mathcal H^{-4}(\mathcal P_3)$ on the ordinary atlas.
These sections are equivariant under positive scaling of the matrices.
Since $\kappa_3$ is nonzero at $x$, its origin stalk is nonzero:
a section with zero origin germ vanishes near the origin, and scaling
would then make it vanish everywhere.
Lemma~\ref{sch:lem:rank-three-bm} makes that scalar group one-dimensional.
Write $\kappa_2(1)=c\kappa_3(1)$ at the origin.
The difference section has zero origin germ, hence vanishes everywhere
by the same argument. Its value at $x$ is $-2c w_2$.
Lemma~\ref{sch:lem:rank-two-decoration} forces $c=0$.
\end{proof}

The incoming filtration admits a specified geometric lift.
Define $\tau_2:\mathbb Q_{\mathfrak M_2}[2]\to\mathcal P_2$
by the complete-flag construction with the first Chern class of
its quotient line inserted. Its scalar degree-$-2$ class is $w_2$.
On the scalar Grassmannian $G=\operatorname{Gr}(2,3)$, let
$H=c_1(\mathbb C^3/U)$, with $\int_G H^2=1$.
Define two classes in $H^{-4}(D_s)$ by
\[
 a=H\,(v_2\otimes v_1),\qquad
 b=\left(Rb_{2,1*}a_{2,1}^*(\tau_2\boxtimes\sigma_1)
                    \circ\eta_{2,1}[4]\right)_s(1).
\]
The first class lies in the degree-two flag contribution. The second
is an actual derived coefficient class, not a chosen splitting of
the cohomology family.

\begin{theorem}[The degree-minus-four scalar image]
\label{sch:thm:degree-four-image}
The classes $a,b$ form a basis of the two-dimensional scalar source
$H^{-4}(D_s)$. Put $w_3=\frac12\kappa_3(1)$ at the scalar point.
Then $w_3$ generates $H^{-4}(P_{3,0})$, and the actual Hall map is
\[
 H^{-4}(\mu_{2,1})(a)=w_3,
 \qquad H^{-4}(\mu_{2,1})(b)=0.
\]
Thus the degree-$-4$ image is the whole target and its kernel is
$\mathbb Q b$. Under the proper Milnor-fibre identification
\eqref{sch:eq:degree-four-gysin}, the Gysin map has rank one.
\end{theorem}

\begin{proof}
The first class generates the first term of the filtration in
Proposition~\ref{sch:prop:degree-four}.
The edge image of $b$ in its second quotient is $w_2\otimes v_1$,
because the scalar stalk of $\tau_2$ is $w_2$ at every plane.
This is nonzero by Lemma~\ref{sch:lem:rank-two-decoration}.
Hence $a,b$ form a basis.

Refine the plane by its complete flag. Inserting its quotient-line
Chern class gives $\ell_2$ on the complete flag of dimension three.
The proper refinement and support naturality used in
Theorem~\ref{sch:thm:images} commute with cup product by this
pulled-back Chern class. Consequently
$\mu_{2,1}(b)=\kappa_2(1)$.
Inserting the outer quotient-line class instead gives $\ell_3$.
The normalized plane input is $\sigma_2=h_2/2$.
Thus $\mu_{2,1}(a)=\kappa_3(1)/2=w_3$.
Lemma~\ref{sch:lem:three-decorations} proves the displayed values
and the nonvanishing of $w_3$.
The target dimension follows from
Lemma~\ref{sch:lem:rank-three-bm}.
\end{proof}

The normalization of $w_3$ is geometric. The section
$\kappa_3/2$ specializes to $w_2$ when a scalar plane is separated
from a line. The class $w_2$ is the positive flag decoration followed
by the determinant-incidence support map. This fixes the generator
without identifying it with an unspecified Borel--Moore basis.

The computation determines one higher scalar degree for the
rank-$(2,1)$ Hall map. The remaining degrees and the higher scalar
images for ranks four and five retain their separate support and
extension calculations.
